\section{Processor Design}

\subsection{Von-Neumann Architecture}

\begin{definition}
The \textbf{Von Neumann model}, proposed in 1946, defines a universal architecture for general-purpose computers consisting of five fundamental components:
\begin{itemize}
    \item \textbf{Memory:} Stores both the program (instructions) and the data.
    \item \textbf{Processing Unit:} Performs arithmetic and logical operations. It contains the Arithmetic Logic Unit (ALU) and small temporary storage (registers).
    \item \textbf{Control Unit:} Orchestrates the execution of instructions. It contains the Program Counter (PC), which tracks the address of the next instruction, and the Instruction Register (IR), which holds the current instruction.
    \item \textbf{Input:} Devices to get information into the computer (e.g., keyboard).
    \item \textbf{Output:} Devices to get information out (e.g., monitor).
\end{itemize}
\end{definition}

\begin{important}
Modern computing is based on two key properties of this model:
\begin{itemize}
    \item \textbf{Stored Program:} Instructions are stored in a linear memory array alongside data. The hardware interprets a memory value as an instruction or data depending on which phase of the cycle it is in.
    \item \textbf{Sequential Processing:} Instructions are processed one at a time in sequence, unless a control flow instruction explicitly changes the Program Counter.
\end{itemize}
\end{important}

\begin{remark}
The process of executing a single instruction follows a specific sequence of phases:
\begin{itemize}
    \item 1. \textbf{FETCH:} Obtain instruction from memory into the IR and increment the PC.
    \item 2. \textbf{DECODE:} Identify the opcode and determine the control signals needed.
    \item 3. \textbf{EVALUATE ADDRESS:} Calculate the memory address of operands if needed (e.g., for loads).
    \item 4. \textbf{FETCH OPERANDS:} Get the required source values from memory or registers.
    \item 5. \textbf{EXECUTE:} Perform the actual operation in the ALU.
    \item 6. \textbf{STORE RESULT:} Write the final value into the destination register or memory.
\end{itemize}
\end{remark}

\begin{definition}
\textbf{Address Space:} The total number of uniquely identifiable memory locations. (LC-3: $2^{16}$; MIPS: $2^{32}$).
\end{definition}

\begin{definition}
\textbf{Addressability:} The number of bits stored at each address (e.g., byte-addressable (8 bit) vs. word-addressable (32 bit)).
\end{definition}

\begin{definition}
\textbf{Endianness:} The order of bytes in a word. \textbf{Big Endian} stores the Most Significant Byte (MSB) at the lowest address, while \textbf{Little Endian} stores the Least Significant Byte (LSB) at the lowest address.
\end{definition}

\begin{example}
To read from memory:
\begin{enumerate}
    \item Load the \textbf{Memory Address Register (MAR)} with the target address.
    \item The memory hardware places the data from that address into the \textbf{Memory Data Register (MDR)}.
\end{enumerate}
To write, both MAR (address) and MDR (data) are loaded before activating the write-enable signal.
\end{example}

\subsection{Instruction Set Architecture (ISA)}

\begin{definition}
The Instruction Set Architecture (ISA) serves as the "contract" or interface between the software (compiler/programmer) and the hardware (microarchitecture). It specifies:
\begin{itemize}
    \item \textbf{Memory Organization:} Address space and addressability.
    \item \textbf{Register Set:} Number, size, and usage of registers (e.g., 8 GPRs in LC-3, 32 in MIPS).
    \item \textbf{Instruction Set:} Opcodes, data types, and addressing modes.
\end{itemize}
\end{definition}

\begin{important}
Instructions are encoded as binary "words" containing:
\begin{itemize}
    \item \textbf{Opcode:} Specifies the operation to perform (e.g., ADD, LD).
    \item \textbf{Operands:} Specify where to find the source data and where to put the result.
\end{itemize}
\end{important}

\begin{lemma}
The semantic gap refers to the distance between high-level language constructs and the ISA. 
\begin{itemize}
    \item \textbf{CISC (Complex Instruction Set Computing):} Small semantic gap; instructions do a lot of work (e.g., x86).
    \item \textbf{RISC (Reduced Instruction Set Computing):} Large semantic gap; simple instructions serve as primitives (e.g., MIPS, LC-3).
\end{itemize}
\end{lemma}

\begin{definition}
Mechanisms used to specify where an operand is located:
\begin{itemize}
    \item \textbf{Immediate:} The operand is a constant within the instruction.
    \item \textbf{Register:} The operand is located in a general-purpose register.
    \item \textbf{PC-Relative:} The address is calculated as $PC + offset$.
    \item \textbf{Indirect:} The instruction points to a memory location that contains the \textit{actual} address of the operand.
    \item \textbf{Base + Offset:} The address is calculated as $Register + offset$.
\end{itemize}
\end{definition}

\subsection{Assembly}

\begin{definition}
\textbf{Machine code} is the binary representation (0s and 1s) processed by hardware. \textbf{Assembly language} is a human-readable representation using mnemonics (e.g., `ADD`) and labels for memory locations.
\end{definition}

\begin{example}
High-level: $a = b + c$
\begin{itemize}
    \item \textbf{LC-3 Assembly:} `ADD R0, R1, R2` (R0 is destination).
    \item \textbf{MIPS Assembly:} `add \$s0, \$s1, \$s2` (\$s0 is destination).
\end{itemize}
\end{example}

\begin{important}
In RISC (Reduced Instruction Set Computing) architectures, operands must be moved into registers before processing.
\begin{itemize}
    \item \textbf{Load (LD/LDR/lw):} Reads data from memory into a register.
    \item \textbf{Store (ST/STR/sw):} Writes data from a register into memory.
\end{itemize}
\end{important}

\begin{remark}
The sequential execution order can be altered using:
\begin{itemize}
    \item \textbf{Conditional Branches:} Executed only if a condition is met. LC-3 uses \textbf{Condition Codes} (N, Z, P) set by previous instructions. MIPS often compares two registers directly (e.g., `beq $s0, $s1, Label`).
    \item \textbf{Unconditional Jumps:} Always change the PC (e.g., `JMP` in LC-3, `j` in MIPS).
\end{itemize}
\end{remark}

\subsection{Microarchitecture}

\begin{definition}
\textbf{Microarchitecture} is the specific implementation of an Instruction Set Architecture (ISA) under given design constraints and goals. It describes how hardware structures, such as registers, ALUs, and buses, are organized to execute the instructions defined by the ISA.
\end{definition}

\begin{remark}
While the ISA acts as the functional contract visible to the software, the microarchitecture is the internal logic that performs the work.
\end{remark}

\begin{important}
A single ISA (such as x86, ARM, or MIPS) \textbf{can be implemented by many different microarchitectures}. This allows manufacturers to create different chips that run the same software but vary in performance, power consumption, or cost. The software remains unaware of whether the hardware uses a single-cycle, multi-cycle, or out-of-order execution strategy.
\end{important}

\begin{remark}
Microarchitects evaluate their designs based on several critical metrics:
\begin{itemize}
    \item \textbf{Performance:} Usually measured by execution time or throughput.
    \item \textbf{Cost:} Often determined by the silicon die area and circuit complexity.
    \item \textbf{Power/Energy:} Critical for thermal management and battery-operated devices.
    \item \textbf{Reliability:} The ability of the system to operate correctly despite aging or environmental interference.
\end{itemize}
\end{remark}

\subsubsection{Single-Cycle Microarchitecture}

\begin{definition}
In a single-cycle microarchitecture, every instruction is executed in exactly one clock cycle. This requires all phases of the instruction cycle, Fetch, Decode, Execute, Memory Access, and Write-back, to be completed within a single period of the global clock.
\end{definition}

\begin{important}
The timing and performance of a single-cycle processor are characterized by:
\begin{itemize}
    \item \textbf{CPI (Cycles Per Instruction):} The CPI is always 1.
    \item \textbf{Clock Cycle Time ($T_c$):} The clock period must be long enough to allow the \textbf{slowest} instruction in the ISA to complete. 
\end{itemize}
\end{important}

\begin{remark}
In most RISC architectures, the lw (load word) instruction defines the critical path. It requires an instruction fetch, a register read, an ALU operation to calculate the address, a data memory read, and a register write. Because the clock must accommodate this long path, faster instructions like add or jump waste a significant portion of the clock cycle waiting for the next edge.
\end{remark}

\subsubsection{Multi-Cycle Microarchitecture}

\begin{definition}
A multi-cycle microarchitecture breaks the execution of a single instruction into multiple, shorter clock cycles. Each instruction takes a variable number of cycles depending on its specific requirements (e.g., an add might take 3 cycles, while a lw takes 5).
\end{definition}

\begin{important}
Multi-cycle designs allow for significant hardware reuse and require internal storage:
\begin{itemize}
    \item \textbf{Functional Unit Reuse:} A single ALU can be used in different cycles of the same instruction (e.g., once to increment the PC and later to perform the arithmetic).
    \item \textbf{Microarchitectural Registers:} Since an instruction spans multiple cycles, intermediate results must be saved in registers that are not visible to the ISA, such as the Instruction Register (IR), Memory Data Register (MDR), and ALUOut.
\end{itemize}
\end{important}

\begin{example}
The phases of a multi-cycle instruction typically include:
\begin{itemize}
    \item \textbf{Cycle 1 (Fetch):} Instruction is fetched and PC is incremented.
    \item \textbf{Cycle 2 (Decode):} Opcode is identified and registers are read.
    \item \textbf{Cycle 3 (Execute):} ALU performs the calculation or address computation.
    \item \textbf{Cycle 4 (Memory):} Data memory is accessed (only for loads/stores).
    \item \textbf{Cycle 5 (Write-back):} Result is written back to the register file.
\end{itemize}
\end{example}

\begin{theorem}
The execution time of a program is defined by the following equation:

$$\text{Time} = \text{Instruction Count} \times \text{Average CPI} \times \text{Clock Cycle Time}$$

While a multi-cycle design has a higher average CPI than a single-cycle design, it can achieve better performance because the clock cycle time ($T_c$) is much shorter, as it only needs to cover the delay of the single longest \textit{stage} rather than the entire instruction.
\end{theorem}

\begin{remark}
The \textbf{control unit} in a multi-cycle processor is implemented as a Finite State Machine (FSM). It must keep track of the current step of an instruction to ensure the correct control signals (like ALUSel or MemWrite) are asserted at the correct time.
\end{remark}

\begin{lemma}
The primary \textbf{downside} of multi-cycle designs is the \textbf{"sequencing overhead"}. Each cycle transition adds a delay due to register setup times and clock-to-Q propagation. If instructions are broken into too many cycles, this overhead can eventually negate the performance benefits of a high clock frequency.
\end{lemma}